\section{Introduction}
\label{sec:introduction}

Evaluating application performance on novel architectural ideas is key to computer architecture research.
However, especially for hardware which does not yet exist, evaluating application performance can be very time-consuming.
For instance, when using cycle-level software simulation, there can be over 100,000$\times$ slowdown.
Even using FPGA-accelerated simulation is significantly slower than \revision{wall-clock} time~\cite{firesim}.
At the same time, modern applications are growing in their computation demands from exascale applications~\cite{exascale} to machine learning training (GPT-4 required 2,150,000 Exa- or $2.15 \times 10^{25}$ operations to complete~\cite{epoch2024trackinglargescaleaimodels}).

To facilitate the evaluation of large applications, there has been a long history of reducing applications to smaller benchmarks, kernels, or short representative intervals.
For instance, Livermore Loops provide a set of microbenchmarks that represent key computational kernels from scientific applications~\cite{livermore_loops}.
It is also common to pull out kernels from an application or to create ``proxy applications''~\cite{exascale-proxy-apps}.
While these methods can reduce very large applications to much more manageable benchmarks, they require application experts to manually identify the important parts of applications and create new programs.

Another approach to reducing the size of applications is to \emph{sample} the application to automatically find \emph{representative intervals} of different application phases.
Many techniques (e.g., SimPoint~\cite{simpoint}, BarrierPoint~\cite{barrier_point}, and LoopPoint~\cite{looppoint}) follow this approach.

In this study, we focus on targeted sampling techniques, which use profiling to automatically select samples.

Prior targeted sampling techniques have three main drawbacks that limit their applicability to modern, realistically sized applications:

\begin{itemize}
    \item \textbf{Samples are expensive to find:} Targeted sampling techniques require running the application from beginning to end once (typically via simulation or slow emulation) to define intervals, analyze program phases, and identify representative samples.
    \item \textbf{Samples are tied to a single executable binary:} These techniques generate samples that are specific to a particular binary, so any binary change necessitates repeating the entire sampling process.
    \item \textbf{Validating the selected samples is infeasible:} Targeted sampling techniques require running the application in detailed simulation to verify that the selected samples accurately represent the program.
\end{itemize}

This study takes a two-part approach.
First, because existing sampling methodologies have mostly been validated in user-mode simulation, we implement the state-of-the-art LoopPoint methodology in gem5's full-system mode and evaluate it with the NPB benchmark suite.
This investigation confirms that LoopPoint can be extended to full-system simulation, but also highlights that the analysis remains time-consuming and that validating sampling accuracy still requires a full baseline simulation---the very bottleneck that sampling aims to avoid.

Motivated by this observation, we present \emph{Nugget}, a framework for agile targeted sampling development.
Because modern applications exhibit varied characteristics, we believe there is no one-size-fits-all sampling methodology or selection algorithm~\cite{survey_phase_classification,accl_warmup_sampled_simulation,characterizing_and_comparing_prevailing_simulation_techniques,analysis_of_statistical_sampling_in_microarchitecture_simulation}.
Therefore, Nugget is designed to be flexible and extensible, allowing researchers to create new sample selection algorithms tailored to their targeted workload and input without the current limitations of targeted sampling techniques.
Nugget's cross-platform compatibility and ability to run on real hardware enable richer insights into both the selected samples and the systems that run them.

The main contributions of this study are:

\begin{itemize}
    \item \textbf{A full-system LoopPoint implementation in gem5}, providing the first evaluation of LoopPoint with the NPB benchmark suite in full-system simulation, establishing a simulation-based baseline for sampling methodology validation.
    \item \textbf{Efficient and portable interval analysis} on real hardware using LLVM, unconstrained by the hardware's microarchitecture or ISA.
    \item \textbf{An LLVM IR-level unit of work}, enabling sample selection algorithms to create and locate samples for a program in a cross-platform manner, agnostic of machine-level instructions, such as architecture-specific instructions and optimizations.
    \item \textbf{A sample creation methodology} that allows samples to run on any platform (e.g., real hardware or simulators).
    \item \textbf{An efficient validation methodology} \revision{that validates the selected samples for the target workload and input on native hardware, integrated into the development workflow.}
\end{itemize}

